\documentclass[11pt]{article}
\usepackage[a4paper,margin=25mm]{geometry}
\usepackage{amsmath,amssymb,mathtools}
\usepackage{fontspec,unicode-math}
\setmainfont{Latin Modern Roman}
\setsansfont{Latin Modern Sans}
\setmonofont{Latin Modern Mono}
\setmathfont{Latin Modern Math}
\usepackage{xurl,hyperref}
\hypersetup{colorlinks=true,urlcolor=blue,linkcolor=blue}
\newcommand{\sref}[2]{\hyperlink{source:#1:#2}{[S#2]}}
\newcommand{\cataloguescope}[1]{\paragraph{Mathematical question.}#1}
\begin{document}
% UnsolvedMath ID 8700059; source catalogue code AMR-086-0059
\setcounter{section}{8700058}
\section{Question 4.17 — Mazur}\label{8700059}
{\small\sffamily UnsolvedMath ID 8700059 · Graph Theory}
\subsection{Definitions and mathematical statement}

\paragraph{Shared notation.}$\mathbb N=\{1,2,\ldots\}$, and $\mathbb Z,\mathbb Q,\mathbb R,\mathbb C$ denote the integers, rationals, reals and complex numbers. Any other conventions are specified locally.


Let $V$ be a smooth algebraic variety over $\mathbb Q$. Give $V(\mathbb R)$ its usual real topology and set $Z=\overline{V(\mathbb Q)}^{\,V(\mathbb R)}$. The condition that $V(\mathbb Q)$ is Zariski dense means it is not contained in a proper Zariski-closed subset of $V$.
\paragraph{Statement recorded in the supplied source.}
If $V(\mathbb Q)$ is Zariski dense in $V$, must $Z$ be a union of connected components of $V(\mathbb R)$? \sref{8700059}{2}
The cited survey records a negative answer to this general question: Colliot-Thélène, Skorobogatov and Swinnerton-Dyer constructed a smooth surface with Zariski-dense rational points whose real closure is not a union of components. The survey subsequently distinguishes the separate abelian-variety case. \sref{8700059}{2}
\subsection{Short English statement}
If rational points are algebraically dense, must their real closure consist of whole real connected components?
\subsection{Sources}
\begin{description}
\item[\hypertarget{source:8700059:1}{S1}] ULAM.AI and the UnsolvedMath Contributors, UnsolvedMath: A Curated Collection of Open Mathematics Problems, record 8700059 (AMR-086-0059). CC BY 4.0. \url{https://huggingface.co/datasets/ulamai/UnsolvedMath}.
\item[\hypertarget{source:8700059:2}{S2}] Michel Waldschmidt, \emph{Open Diophantine Problems}, arXiv:math/0312440v2 (24 January 2004), Question 4.17 and following discussion, p.42. \url{https://arxiv.org/abs/math/0312440v2}.
\end{description}
\label{end:8700059}
\end{document}
